\chapter{Ei eiga kunnskapsform}

\begin{motto}
Det finst ikkje éin måte å vite på.\\ Det finst minst tre, og design er den tredje.
\end{motto}

Når \textsc{Grunnlaust} hevdar at design manglar ei eiga kunnskapsform, er det ikkje ein observasjon. Det er ein dom som føreset svaret. For å sjå at design ikkje veit noko på sin eigen måte, må ein alt ha bestemt at den einaste måten å vite på er den vitskapen brukar — å skildre verda som ho er, i påstandar som let seg etterprøve. Mot den målestokken kjem design til kort, av di design ikkje held på med å skildre verda som ho er. Men då er ikkje konklusjonen at design er utan kunnskap. Konklusjonen er at vi har brukt feil målestokk. Dette kapitlet legg fram den positive påstanden boka kviler på: at design har eit eige studieobjekt, eigne metodar og eigne verdiar, at desse er kartlagde i ein førti år gammal forskingslitteratur, og at dei utgjer ein heil kunnskapskultur — ikkje eit tomrom kledd i slagord.

Eg vil vere tydeleg på kva slags påstand dette er, for å skilje han frå ein svakare. Det svake forsvaret seier: design har \emph{òg} ei form for kunnskap, ved sida av den eigentlege, vitskaplege. Det er eit dårleg forsvar, av di det godtek at den vitskaplege kunnskapen er den eigentlege og tildeler design ein andreplass. Det sterke forsvaret, som er det boka fører, seier noko anna: at det finst fleire jamstilte måtar å vite på, kvar med sitt objekt og sine kriterium, og at design er ein av dei i full forstand — ikkje ein blekare slektning av vitskapen, men eit sidestykke til han. Skilnaden er ikkje retorisk. Han avgjer kven som ber bevisbyrda. Om design berre er ei underordna form for kunnskap, må det rettferdiggjere seg framfor vitskapen. Om det er ei jamstilt form, er det kritikaren som må vise kvifor den eine forma skulle vere den einaste lovlege. Resten av kapitlet held fram med den sterke påstanden, fordi det er den som faktisk er sann.

\section{Det Cross faktisk hevda}

Frasen «designerly ways of knowing» er blitt så mykje sitert at ho lett kan klinge som ei besverjing — ein måte å hevde at design veit noko utan å seie kva. Det er verdt å gå tilbake til kva Nigel Cross faktisk hevda, fordi påstanden hans er presis og innhaldsrik. I den korte programartikkelen frå 1982 \autocite{cross1982}, og i den seinare samlinga som ber same tittel \autocite{cross2006}, set Cross fram ein tese om at det finst tre breie kulturar for menneskeleg kunnskap, ikkje éin. Den fyrste er vitskapen, hvis objekt er den naturlege verda, hvis metode er det kontrollerte eksperimentet og analysen, og hvis verdiar er objektivitet, rasjonalitet og søken etter sanning. Den andre er humaniora, hvis objekt er den menneskelege røynsla, hvis metode er analogi, metafor og kritisk tolking, og hvis verdiar er subjektivitet, rettferd og søken etter meining. Den tredje er design — eller meir vidt, det Cross etter Simon kallar det kunstige, det menneskeskapte — hvis objekt er den materielle verda slik ho \emph{kunne vore}, hvis metode er modellering, mønsterdanning og syntese, og hvis verdiar er praktisk dugleik, oppfinnsemd og omsut for det høvelege.

Det avgjerande er at dette ikkje er ei rangering. Cross hevdar ikkje at design er ein ufullkomen vitskap, ein vitskap som enno ikkje har vakse opp. Han hevdar at design er ein annan ting, med eit anna objekt og difor med ein annan logikk. Vitskapen spør kva som er tilfellet; design spør kva som bør lagast. Dei to spørsmåla har ulik grammatikk. Frå ei skildring av korleis verda er, følgjer aldri åleine kva ein bør gjere med henne — dette er eit gammalt poeng i filosofien, og det rammar rett i kjernen av kravet \textsc{Grunnlaust} stiller. Å be design om å grunngje vala sine slik vitskapen grunngjev påstandane sine, er å be om at ein skal utleie eit \emph{bør} av eit \emph{er}. Det lèt seg ikkje gjere, ikkje fordi design er slapp, men fordi det er to ulike slag arbeid.

Det er verdt å dvele ved det innhaldet Cross faktisk fyller den tredje kulturen med, for det er her påstanden går frå program til substans. Cross peikar på at designarar gjennom utdanninga si tileignar seg ein eigen kompetanse i å handtere det han kallar dårleg definerte problem — oppgåver utan ein fasit og utan ein uttømmande spesifikasjon — ved hjelp av løysingsfokuserte strategiar heller enn problemfokuserte. Dei lærer å bruke ikkje-verbale, grafiske og romlege medium for å tenkje med: skissa, modellen, diagrammet er ikkje berre måtar å vise fram ein ferdig idé på, men sjølve reiskapane resonnementet skjer i. Dei lærer å lese eit objekt eller ein situasjon og slutte seg til kva det \emph{kunne} bli — ein form for syntese der mange uforeinlege krav blir samla i ein einskild, konkret ting. Dette er ikkje vage dygder. Det er namngjevne, studerte kompetansar, kartlagde gjennom protokollstudiar og observasjon av korleis designarar arbeider, og dei utgjer eit positivt svar på spørsmålet om kva designkunnskap \emph{er}. Når kritikaren spør kva som er innhaldet bak frasen, er svaret altså ikkje eit forlegent fråvær, men ein katalog: ein måte å ramme problem på, ein måte å tenkje gjennom medium på, og ein måte å syntetisere mot det høvelege.

Det same poenget viser seg i Cross sitt seinare skifte i ordbruk, frå \emph{designvitskap} til \emph{designdisiplin} \autocite{cross2001discipline}. Dette skiftet er ofte blitt lese som ei innrømming: design klarte ikkje å bli ein vitskap, så ein nøgde seg med å kalle det ein disiplin. Lesinga er feil. Skiftet er ikkje eit nederlag, men ei korrigering. Herbert Simon hadde drøymt om ein eksakt vitskap om det kunstige, modellert tett på naturvitskapen \autocite{simon1969sciences}; Cross sitt poeng er at denne draumen bad design om å bli noko anna enn det er. Ein disiplin er ikkje ein mislukka vitskap. Det er ei form for organisert kunnskap med eige objekt, eigne metodar og eigne kriterium for kva som tel som godt arbeid. Jussen er ein disiplin og ingen vitskap; det same er historiefaget og medisinen i si kliniske form. Å kalle design ein disiplin er ikkje å seie at det er mindre enn ein vitskap. Det er å seie kva slags ting det er.

\section{Archer og fødselen av eit fagfellevurdert felt}

Det kan innvendast at alt dette er programmatikk — fine ord om kva design \emph{kunne} vere, utan den institusjonelle ryggrada som skil ein disiplin frå ein ambisjon. Men ryggrada finst, og ho har ein datert opphavsstad. I 1979, i den fyrste årgangen av tidsskriftet \emph{Design Studies}, sette Bruce Archer fram tesen at design er ein tredje kunnskapsområde ved sida av vitskapane og humaniora, med eit eige objekt for systematisk gransking \autocite{archer1979discipline}. Det som gjer Archer sitt bidrag avgjerande, er ikkje berre påstanden, men forma han kom i. Archer skreiv ikkje eit manifest. Han skreiv ein fagfellevurdert artikkel, i eit fagfellevurdert tidsskrift, med ein redaksjon, ein referansepraksis og ein lesarkrins av andre forskarar som kunne motseie han. Det er sjølve den handlinga — å føre påstanden inn i eit ordskifte med reglar — som grunnlegg ein disiplin.

Her ligg eit svar på \textsc{Grunnlaust} sitt bilete av design som eit handverk forkledd som fag. Eit handverk har meister og lærling, men det har ikkje eit fagfellevurdert ordskifte der påstandar blir prøvde av jamlikar som ikkje står i lærdomsforhold til kvarandre. Frå Archer og framover har design hatt nett det. \emph{Design Studies} kom i 1979; \emph{Design Issues} følgde i 1984; sidan har det vakse fram ei rekkje fagfellevurderte tidsskrift, internasjonale konferansar med fagfellevald program, og doktorgradsprogram med opponentar og avhandlingar. Eit felt med dette apparatet er per definisjon ikkje berre eit handverk. Det er eit handverk \emph{og} ein disiplin — praksisen og den organiserte refleksjonen over praksisen, slik medisinen er både legekunst og medisinsk vitskap. At designdisiplinen er ung — knapt eit halvt hundreår gammal i fagfellevurdert form — gjer han ikkje uekte. Det gjer han ung.

Archer hadde dessutan gjort noko meir enn å skrive ein artikkel; han hadde, gjennom arbeidet sitt ved Royal College of Art, bygd opp eit forskingsmiljø som tok mål av seg å studere design systematisk. Det vesentlege i forma hans var insisteringa på at design er noko ein kan forske \emph{på} og forske \emph{gjennom} — at det finst eit kunnskapsobjekt her som ikkje er redusert til verken kunsthistorie eller ingeniørvitskap, men som krev sine eigne omgrep og sine eigne metodar. Denne tanken, at design er eit tredje hovudområde i allmennutdanninga jamsides vitskapen og humaniora, var radikal då han kom, og ho la grunnlaget for ein heil generasjon designforsking. Poenget mot \textsc{Grunnlaust} er presist: ein disiplin blir ikkje til ved at nokon utropar henne, men ved at nokon byrjar å føre eit dokumentert, fagfellevurdert ordskifte om eit avgrensa kunnskapsobjekt. Det er nett det som skjedde, og det skjedde på ein datert stad.

Ein ærleg merknad høyrer med. Eit fagfellevurdert apparat garanterer ikkje at kunnskapen som flyt gjennom det er kumulativ, samanhengande eller einig med seg sjølv; det garanterer berre at han er prøvd. Spørsmålet om kumulasjon kjem seinare i boka, og det er eit reelt spørsmål. Men poenget her er smalare og held: påstanden om at design ikkje har ei disiplinær form er empirisk feil. Forma har vore der, datert og dokumentert, sidan 1979.

\section{Nelson og Stolterman: eit heilt fundament, ikkje eit fråvær}

Den hardaste versjonen av skuldinga er ikkje at design manglar eit apparat, men at det manglar eit fundament — ein eksplisitt redegjord teori om kva designkunnskap \emph{er}, korleis han skil seg frå anna kunnskap, og kva som gjer han gyldig. Mot denne skuldinga finst det eit verk som er vanskeleg å koma forbi: \emph{The Design Way} av Harold Nelson og Erik Stolterman \autocite{nelson2012designway}. Boka gjer presis det \textsc{Grunnlaust} seier ikkje finst. Ho legg fram, eksplisitt og samanhengande, ein metafysikk og ein epistemologi for design forstått som ein eigen «kultur av gransking og handling» ved sida av vitskapen og kunsten.

Kjernen i Nelson og Stolterman sitt argument er eit skilje mellom det sanne, det ideelle og det røynlege. Vitskapen søkjer det sanne — det som er tilfellet, uavhengig av oss. Kunsten og etikken kan søkje det ideelle — det fullkomne, det som bør vere. Design søkjer noko tredje, som dei kallar det røynlege: den \emph{partikulære}, ikkje-naudsynte tingen som blir til når nokon vel å la henne bli til, her, no, under desse vilkåra, for desse menneska. Designkunnskap er kunnskap om korleis ein bringar det røynlege til verda — ikkje kunnskap om allmenne lover, men om det einskilde tilfellet og om det høvelege grepet i det. Dei kallar denne dugleiken \emph{design judgement}, designskjøn: den samansette evna til å vege uforeinlege omsyn mot kvarandre og koma fram til eit heilt, forsvarleg grep som ingen formel kunne ha rekna ut. Skjønet er ikkje irrasjonelt. Det er ein eigen form for rasjonalitet, retta mot det partikulære heller enn det allmenne.

Nelson og Stolterman går vidare enn å skildre objektet; dei spesifiserer dei intellektuelle dygdene design krev. Skjønet bryt dei opp i fleire slag — eit skjøn for kva som er sant i situasjonen, eit for kva som har verdi, eit for kva som høver, eit for korleis delane skal komponerast til ein heilskap, og eit avgjerande skjøn for kva som er det rette grepet \emph{her}, i dette eine tilfellet. Dei løftar fram \emph{komposisjon} som ein eigen kunnskapsmodus: evna til å setje saman element som ikkje følgjer av kvarandre logisk, slik at heilskapen blir meir enn summen. Og dei plasserer \emph{omsut} og \emph{tenesta} i sentrum — design er for dei ikkje ein nøytral teknikk, men ein intensjonell handling på vegner av nokon, retta mot å gjere ein tilstand betre for nokon. Dette er ein eksplisitt etikk innebygd i sjølve kunnskapsforma, ikkje bolta på etterpå. Heile bygnaden er meir enn rik nok til å motseie påstanden om eit tomrom: her er ein ontologi (læra om det røynlege), ein epistemologi (læra om designskjønet) og ein verditeori (læra om omsut og tenesta), alt utleidd frå design sjølv og ikkje lånt frå eit anna fag.

Det viktige for forsvaret er at dette er positivt innhald. Det er ikkje ein klage over at vitskapen ikkje passar design; det er ein utbygd teori om kva designkunnskap er i seg sjølv, med eigne omgrep — det røynlege, designskjønet, komposisjon, omsut, det ultimate partikulære — som ein kan lære, prøve og strides om. Når kritikaren seier at design «berre» har slagord, har han ikkje lese den litteraturen som leverer det motsette: ein eksplisitt redegjord tredje kunnskapskultur. Ein kan vere usamd med Nelson og Stolterman; ein kan finne teorien deira for laus på enkelte punkt. Men ein kan ikkje seie at fundamentet manglar utan å late som om boka ikkje finst. Eit fundament som er framlagt og som ein vel å ikkje sjå, er framleis eit fundament.

\section{Kvifor dette ikkje er ein retrettposisjon}

Ein skarp lesar vil her kjenne ein mistanke, og han fortener eit svar. Er ikkje heile dette grepet ein retrett? Fyrst krev design å vere ein vitskap; så, når det ikkje held, omdøyper det seg til ein «eigen kunnskapskultur» som per definisjon ikkje kan målast mot vitskapen og difor ikkje kan motseiast. Er ikkje «designerly ways of knowing» berre ein måte å gjere seg uangripeleg på?

Innvendinga er god, og ho må takast på alvor, for ho råkar ein reell freistnad i faget. Men ho råkar ikkje argumentet her, av to grunnar. For det fyrste er ikkje den tredje kunnskapskulturen ein kategori design fann opp for å sleppe unna. Skiljet mellom å vite \emph{at} og å vite \emph{korleis}, mellom kunnskap om det naudsynte og kunnskap om det partikulære, er eldgammalt og langt eldre enn designforskinga — det går attende til Aristoteles' skilje mellom \emph{episteme}, \emph{techne} og \emph{phronesis}. Design plasserer seg i ein kategori som var der frå før. For det andre, og viktigare: den tredje kulturen er ikkje uangripeleg. Han har sine eigne kriterium for godt og dårleg arbeid, og dei er strenge. Eit designgrep kan vere klønete, uærleg, dårleg gjennomtenkt, blindt for verknadene sine; ein designteori kan vere upresis, sjølvmotseiande eller tom. Schön viste at den indeterminerte praksisen er full av dommar som kan vere rette eller gale \autocite{schon1983}, og det er nett det neste kapitlet skal vise i detalj. Å ha eigne kriterium er ikkje det same som å vere fri frå kriterium. Det er berre å nekte å låne kriteria til eit anna fag.

Skilnaden mellom ein ekte tredje kultur og ein retrett ligg her: retretten gjer seg umogleg å ta feil, medan den ekte kulturen berre flyttar spørsmålet om feil over på sine eigne premissar. Design kan ta feil — ikkje ved å spå gale om naturen, men ved å lage gale ting for gale grunnar. Det er ein fullverdig måte å kunne ta feil på, og difor ein fullverdig måte å vite på.

\vignett

Men om designkunnskap i så stor grad er taus og skjønsbasert, melder den hardaste innvendinga seg: er han då i det heile overførbar, eller er han privat — låst inne i den einskilde utøvaren, umogleg å prøve utanfrå? Det er nett der neste kapittel tek til.

\vidarelesing{\fullcite{cross2006}; \fullcite{cross1982}; \fullcite{cross2001discipline}; \fullcite{archer1979discipline}; \fullcite{nelson2012designway}; \fullcite{schon1983}.}
